% ЗӨВХӨН БҮТЭН НООРОГТ (main.tex) хэвлэгдэнэ. Багшид үзүүлэх ноорогт (main_teacher.tex) орохгүй.
% 3.8 бодит асуултаар хийх үнэлгээ. 2026.11.23 (багш): яриа таних аргын сонголт ба 12 бичлэгийн туршилтыг
% Chapter3_Methodology.tex-ийн 3.7 «Яриа таних аргыг сонгох туршилт» руу шилжүүлсэн (багшид үзүүлэх ноорогт).
\section{Бодит асуултаар хийх үнэлгээ}
\label{sec:exp-tree}

Энэ үнэлгээ судалгааны хоёр асуултад (\ref{sec:rq}-р хэсэг) хариулах бөгөөд ресепшнээс бодит асуулт цуглуулсны дараа гүйцэтгэнэ. Түүнээс өмнө системийг \ref{sec:evaluation-method}-р хэсэгт тодорхойлсон, энэ ажлын хүрээнд бэлтгэсэн 40 асуултын багцаар урьдчилан үнэлнэ.
% 2026.11.15 (хэрэглэгч): «хоёр аргыг харьцуулах туршилт» гэж нэрлэхгүй. 2026.11.15 (3): нэмэлт оноогүй
% үеийн тоог суурь болгож нэг хүснэгтэд тайлагнана (хэрэглэгч сонгосон, 5 хянагчийн шалгалтын Ш1).

\textbf{Хайлтын арга.} Үнэлэх хайлт нь BM25 ба E5-ийн жагсаалтыг нэгтгэсэн хайлт (\ref{sec:rag-method}-р хэсэг) бөгөөд түүн дээр мэдээллийн ангиллын дээд түвшний 11 хүрээгээр нэмэлт оноо өгөх тохиргоотой, харин дэд сэдвийн түвшнийг хайлтад ашиглаагүй. Систем асуултын хамаарах хүрээг түлхүүр үгээр тодорхойлох бөгөөд тухайлбал «дадлага» гэсэн үг орсон асуултыг «Карьер, хамтын ажиллагаа» хүрээнд хамааруулна. Дараа нь бүх эх сурвалжаас хайж, тухайн хүрээний эх сурвалжид нэмэлт оноо өгнө (\ref{sec:lit-tree}-р хэсэг). Харин олон сэдэвтэй асуултад хоёр ба түүнээс олон хүрээ тэнцүү таарсан тохиолдолд нэмэлт оноо өгөхгүй. Хариулт олдоогүй үед систем хандах алба, өрөөг заадаг боловч алба, өрөөг том хүрээгээр нь заах нь буруу алба руу чиглүүлэх эрсдэлтэй. Тухайлбал «ажлын байр» гэсэн үг «Карьер, хамтын ажиллагаа» хүрээнд таарах боловч ресепшнд үүнийг сургуульд ажилд орох хүн асуудаг тул алба, өрөөний жагсаалтыг том хүрээгээр биш, ресепшний сэдэв тус бүрээр авна. Сэдэв тодорхойгүй тохиолдолд систем ерөнхий лавлах руу чиглүүлэх бөгөөд энэ алхам жагсаалт батлагдсаны дараа ажиллана (\ref{sec:rag-method}-р хэсэг).

\textbf{Асуулт.} Ресепшн, Facebook-ийн чатботыг хариуцдаг багш, карьерын төвөөс хүмүүсийн бодитоор асуудаг асуултыг цуглуулна. Ресепшний ажилтны өгсөн жагсаалт нь өгүүлбэр биш, сэдэв юм. Иймд ажилтнаас сэдэв бүрд хүмүүсийн асуудаг хэлбэрээр нь асуултыг бичүүлж авна. Эдгээр нь зочдын хэлсэн үгийн бичлэг биш, ажилтнуудын санаж бичсэн асуулт болохыг үр дүнг тайлбарлахдаа харгалзана. Зорилтот тоо нь дор хаяж 40 асуулт, үүнээс дор хаяж 10 нь хариулт нь сургуулийн мэдээлэлд байхгүй асуулт байна. Хөгжүүлэлтийн үед бэлдсэн асуулт болон 11-р сарын 12-нд системийг урьдчилан шалгахад ашигласан жишээ асуултыг энэ үнэлгээнд оруулахгүй.

\textbf{Тэмдэглэгээ.} Асуулт бүрт зөв хүрээ, зөв эх сурвалж, хүлээгдэх үр дүнг (хариулах эсвэл хандах алба, өрөөг заах) системийн хариултыг харахаас өмнө тэмдэглэнэ. Хариулт хэд хэдэн хуудас эсвэл мөрөнд байвал тэдгээрийг бүгдийг зөв эх сурвалж гэж тооцно. Хандах зөв алба, өрөөг системийн жагсаалтаас биш, ресепшний ажилтны хариултаас авна. Тэмдэглэгээг энэ ажлын хүрээнд хийж, асуултын нэг хэсгийг сургуулийн ажилтнаар шалгуулан, зөрсөн тоог тайлагнана.

\textbf{Хуваалт.} Асуултыг цуглуулсан газар (ресепшн, чатбот, карьерын төв) болон хариулттай эсэхээр нь бүлэглэж, бүлэг бүрийг санамсаргүйгээр хоёр тэнцүү хэсэгт хуваана. Сондгой тоотой бүлгийн илүү асуултыг тохируулах хэсэгт оруулна. Хуваасны дараа асуулт, тэмдэглэгээг өөрчлөхгүй бөгөөд хуваасан огноог тэмдэглэнэ. Эхний хэсгээр түлхүүр үг ба нэмэлт онооны хэмжээг тохируулна. Хоёр дахь хэсгийг тохируулахад харахгүй, зөвхөн эцсийн үнэлгээнд ашиглана.

\textbf{Хэмжүүр.} Эцсийн үнэлгээний хэсэгт дараах тоог $k/N$ хэлбэрээр тайлагнана: зөв эх сурвалж эхний 1 ба эхний 3 үр дүнд орсон асуултын тоо; хариулт байхгүй асуултаас зөв алба, өрөөг заасан тоо; хариулттай асуултыг буруугаар алба руу чиглүүлсэн тоо. Эхний тоог мэдлэгийн модны хүрээгээр нэмэлт оноо өгөөгүй ба өгсөн хоёр тохиргоонд зэрэг тайлагнаж, асуулт бүрээр сайжирсан, муудсан, өөрчлөгдөөгүйг тоолно. Нэмэлт оноо өгөөгүй үеийн тоо нь суурь болно. Хүрээг зөв тодорхойлсон асуултын тоог туслах хэмжүүр болгож, буруу тодорхойлсон жишээг хамт тайлагнана. Асуулт цөөн байх тул үр дүнг хувиар биш, тоогоор нь тайлагнана.

\textbf{Шийдвэрийн дүрэм.} Нэмэлт оноотой тохиргоо эцсийн үнэлгээний хэсэгт суурь тооноос илүү гараагүй бол нэмэлт оноог идэвхгүй үлдээнэ. Хэдэн асуултаар илүү гарвал идэвхжүүлэх босгыг асуултыг хуваахтай зэрэг, үр дүнг харахаас өмнө тогтооно. Алба, өрөөний жагсаалт батлагдаагүй бол хоёр дахь судалгааны асуултын тоог тайлагнахгүй бөгөөд үүнийг хязгаарлалтад бичнэ.

\textbf{Хязгаар.} Үнэлгээг бичгэн асуултаар хийнэ. Дуугаар асуусан асуултыг энэ үнэлгээнд тусад нь хэмжихгүй. Зорилтот тоогоор эцсийн үнэлгээний хэсэгт 20 орчим асуулт үлдэнэ. Ийм багцын тоо нь энэ сургуулийн энэ багц асуултад систем хэрхэн ажилласныг л харуулна. Өөр асуулт, өөр сургуульд ерөнхийлөхгүй бөгөөд хоёр тохиргооны нэг, хоёр асуултын зөрүүг ялгаа гэж дүгнэхгүй.
% Асуулт цуглуулсны дараа: асуултын тоо, хуваалт, хуваасан огноог энд нэм.
